% !TEX program = pdflatex
\documentclass[a4paper,12pt]{article}
\usepackage[T5]{fontenc}
\usepackage[utf8]{inputenc}
\usepackage{amssymb,epsfig,latexsym,multicol,array,hhline,fancyhdr}
\usepackage{vntex}
\usepackage{mathptmx}
\usepackage{titlesec}
\usepackage{amsmath}
\usepackage{lastpage}
\usepackage{enumerate}
\usepackage{color}
\usepackage{graphicx}
\usepackage{array}
\usepackage{tabularx, caption}
\usepackage{longtable}
\usepackage{multirow}
\usepackage{multicol}
\usepackage[top=2.5cm, bottom=2.5cm, left=3cm, right=2cm]{geometry}
\usepackage{setspace}
\usepackage{enumitem}
\usepackage{xcolor}
\usepackage[ruled,vlined]{algorithm2e}
\usepackage{hyperref}
\hypersetup{urlcolor=blue,linkcolor=black,citecolor=black,colorlinks=true}
\usepackage{indentfirst}
\usepackage{booktabs}
\usepackage{float}

\usepackage{amsthm}
\newtheorem{theorem}{{\bf Định lý}}
\newtheorem{proposition}{{\bf Mệnh đề}}
\newtheorem{corollary}[proposition]{{\bf Hệ quả}}

\AtBeginDocument{\renewcommand*\contentsname{Mục lục}}
\AtBeginDocument{\renewcommand*\refname{Tài liệu tham khảo}}
\AtBeginDocument{\renewcommand*\figurename{Hình}}
\AtBeginDocument{\renewcommand*\tablename{Bảng}}

\setlength{\headheight}{28pt}
\pagestyle{fancy}
\fancyhead{}
\fancyhead[L]{\small\bf ShortStop --- Phân tích chuyên sâu}
\fancyhead[R]{\small\bf Cong-Tu Le}
\fancyfoot{}
\fancyfoot[L]{\scriptsize Counterexample-Guided Short-Horizon Shielding of Generative Robot Policies}
\fancyfoot[R]{\scriptsize Trang {\thepage}/\pageref{LastPage}}
\renewcommand{\headrulewidth}{0.3pt}
\renewcommand{\footrulewidth}{0.3pt}

\setcounter{secnumdepth}{3}
\setcounter{tocdepth}{3}
\renewcommand{\thesection}{\Roman{section}}
\renewcommand{\thesubsection}{\arabic{section}.\arabic{subsection}}
\renewcommand{\thesubsubsection}{\arabic{section}.\arabic{subsection}.\arabic{subsubsection}}

\setstretch{1.25}
\setlength{\parindent}{1.27cm}
\setlength{\parskip}{6pt}

\titleformat{\section}[block]
{\normalfont\large\bfseries\centering}
{\thesection.}{0.5em}{\MakeUppercase}
\titleformat{\subsection}
{\normalfont\normalsize\bfseries}
{\thesubsection.}{0.5em}{}
\titleformat{\subsubsection}
{\normalfont\normalsize\bfseries\itshape}
{\thesubsubsection.}{0.5em}{}
\titlespacing*{\section}{0pt}{18pt}{9pt}
\titlespacing*{\subsection}{0pt}{12pt}{6pt}
\titlespacing*{\subsubsection}{0pt}{9pt}{4pt}

\let\stdsection\section
\renewcommand{\section}{\clearpage\stdsection}

\newcommand{\st}{\textsc{ShortStop}}

\title{\textbf{Phân tích chuyên sâu: Nền tảng, Cách hiện thực, \\ và Vị trí trong Literature của \st{}} \\[4pt]
\large Counterexample-Guided Short-Horizon Shielding of Generative Robot Policies}
\author{Cong-Tu Le \\ \small Tài liệu nghiên cứu (research note) --- đọc kèm với \emph{ShortStop\_Research\_Plan.tex}}
\date{20/08/2026}

\begin{document}
\maketitle
\tableofcontents

%=====================================================================
\section{Mục tiêu và phạm vi tài liệu}
%=====================================================================
Tài liệu này trả lời sâu 5 câu hỏi về bài báo \st{} (\emph{Counterexample-Guided Short-Horizon Shielding of Generative Robot Policies}, nộp ICRA 2026, double-blind):
\begin{enumerate}[label=(\arabic*)]
  \item \textbf{Nền tảng} bài báo dựa trên những khối kiến thức nào, và bài báo kế thừa/tránh gì từ mỗi khối.
  \item \textbf{Hướng hoạt động} --- insight cốt lõi biến một bài toán bất khả thi (verify hình thức một generative policy) thành một bài toán khả thi (verify hậu quả ngắn hạn của một chunk hành động cụ thể).
  \item \textbf{Cách hiện thực} --- toán học và thuật toán chi tiết của vòng lặp P-R-C-S, cùng bộ ba định lý an toàn.
  \item \textbf{Cái dễ / cái khó} khi hiểu và tái lập bài báo, và \textbf{vì sao đây là một đóng góp mới/khó} chứ không phải một sự ghép nối tầm thường của các kỹ thuật đã biết.
  \item \textbf{4 bài báo liên quan} --- xuất phát từ danh sách reference gốc của \st{} rồi mở rộng ra literature 2024--2025, phân tích điểm mạnh/yếu của từng bài so với \st{}.
\end{enumerate}
Nội dung được tổng hợp trực tiếp từ văn bản đầy đủ của \st{} (\texttt{main (3).txt/pdf}, gồm Sec.~I--IX và Phụ lục A--B), đối chiếu với \emph{ShortStop\_Research\_Plan.tex}, và bổ sung bằng tra cứu học thuật (arXiv/Semantic Scholar) cho phần related-work mở rộng.

%=====================================================================
\section{Bức tranh toàn cảnh cho người ngoài ngành: \\ ``control'' đã đi vào ``policy'' như thế nào}\label{sec:bigpicture}
%=====================================================================
Phần này viết cho người \emph{chưa} ở trong ngành robotics/control --- mục tiêu là hiểu được \st{} \textbf{đang tiếp nối cái gì}, trước khi đi vào chi tiết kỹ thuật ở các phần sau. Nếu bỏ qua phần này, các thuật ngữ như "reachtube", "STL robustness", "shield" sẽ nghe như biệt ngữ rời rạc, không thấy được mạch phát triển liên tục phía sau nó.

\subsection{"Control" là gì, và vì sao nó luôn cần một thứ "phanh lại" có bảo đảm}
Một hệ điều khiển (control system) --- ví dụ máy điều nhiệt, tay lái ô tô, cánh tay robot --- luôn có một vòng lặp: \emph{đo trạng thái hiện tại} $\to$ \emph{quyết định hành động} $\to$ \emph{tác động lên thế giới} $\to$ đo lại. Ngành control theory cổ điển (PID, LQR) tồn tại để trả lời câu hỏi: \emph{làm sao chọn hành động sao cho hệ luôn ở gần mục tiêu và không bao giờ đi vào vùng nguy hiểm, dù có nhiễu?} Điểm mấu chốt của control cổ điển: luật điều khiển được \textbf{viết ra từ phương trình vật lý} (mô hình động lực học đã biết), nên người thiết kế \emph{chứng minh được} bằng toán rằng hệ sẽ ổn định/an toàn --- đây là một \textbf{guarantee}, không phải một lời hứa suông.

\subsection{"Policy" là gì, và vì sao nó dần thay chỗ cho controller viết tay}
Khi hệ quá phức tạp để viết phương trình tay (ví dụ một cánh tay robot 7 khớp cầm vật mềm), người ta chuyển sang \textbf{học} luật điều khiển từ dữ liệu thay vì suy ra từ vật lý --- luật điều khiển học được này gọi là \textbf{policy}: một hàm số (thường là mạng nơ-ron) nhận quan sát (ảnh camera, cảm biến khớp) và trả ra hành động. Hai cách học chính: \emph{reinforcement learning} (thử-sai, tối ưu thưởng) và \emph{imitation learning} (bắt chước dữ liệu demo của con người). Cái được: policy tổng quát hoá tốt hơn cho các tác vụ khó mô hình hoá tay. Cái mất: policy là một \textbf{hộp đen} --- không còn phương trình để chứng minh nó luôn an toàn.

\subsection{Dòng chảy lịch sử: từ luật điều khiển viết tay đến policy sinh (generative policy)}
Bảng dưới đặt \st{} vào đúng vị trí của nó trong một mạch phát triển liên tục, \textbf{không phải một bước nhảy rời rạc}:

\begin{table}[H]
\centering\small
\begin{tabularx}{\textwidth}{@{}l X X X@{}}
\toprule
Giai đoạn & Ai/cái gì quyết định hành động? & Bảo đảm an toàn đến từ đâu? & Cái giá phải trả \\
\midrule
PID / LQR (control cổ điển) & Công thức toán suy từ phương trình vật lý đã biết & Chứng minh giải tích trực tiếp trên phương trình & Cần biết chính xác mô hình vật lý; không tổng quát hoá cho tác vụ phức tạp \\
MPC (\emph{-- đây là nền tảng bạn đã có từ paper trước, xe tự hành --}) & Giải một bài toán tối ưu có ràng buộc ở mỗi bước (receding horizon) & Ràng buộc được đưa trực tiếp vào bài toán tối ưu (QP) & Cần mô hình động lực tractable; chi phí giải QP mỗi bước \\
Reinforcement learning / imitation learning (policy học từ dữ liệu) & Một mạng nơ-ron học từ thử-sai hoặc từ demo & Không có --- policy có thể học một hành vi vi phạm mà không ai biết trước & Mất khả năng chứng minh; phải kiểm định bằng thực nghiệm \\
Safe RL / CBF-QP (cố gắng lấy lại bảo đảm) & Policy học được, nhưng bị một lớp ràng buộc (barrier function) chỉnh lại & Ràng buộc hình học (barrier) áp lên \emph{một} hành động tại một thời điểm & Cần thiết kế tay barrier; chỉ sửa được 1 hành động, không phải cả chuỗi \\
Diffusion / flow-matching policy (\textbf{generative policy}) & Một sampler sinh xác suất, đề xuất cả \emph{chuỗi} hành động tương lai (action chunk) cùng lúc & Không có certificate nào --- chunk được thực thi ngay khi sinh ra & Sampler cỡ trăm triệu tham số, đa modal, hộp đen hoàn toàn \\
\textbf{\st{} (bài báo đang phân tích)} & Vẫn là generative policy ở trên đề xuất, \textbf{nhưng} một lớp kiểm tra độc lập đứng ngoài để lọc/sửa & Quay lại đúng công cụ của control cổ điển (reachable set, invariant set) --- nhưng áp lên \emph{hậu quả} của policy, không áp lên bản thân policy & Phải calibrate model-error, xây fallback controller, chấp nhận guarantee có điều kiện \\
\bottomrule
\end{tabularx}
\caption{Dòng chảy từ control viết tay đến generative policy, và vị trí của \st{}. Đọc theo chiều dọc từ trên xuống là đọc theo thời gian.}
\end{table}

\subsection{Dẫn từng bước vào Diffusion Policy và Flow Policy}\label{sec:diffusion-intro}
Dòng áp chót của bảng trên gói trong hai chữ "diffusion / flow-matching policy" rất nhiều thứ. Phần này dẫn từng bước, không cần biết trước gì về generative model.

\subsubsection{Bước 1 --- vì sao "học một hàm số trả ra 1 hành động" (hồi quy) lại không đủ}
Cách đơn giản nhất để học policy từ dữ liệu demo là hồi quy (regression): huấn luyện một mạng nơ-ron $f(\text{quan sát}) \to \text{hành động}$ bằng cách tối thiểu hoá lỗi bình phương với hành động trong demo. Cách này thất bại khi tác vụ có \textbf{nhiều lời giải đúng như nhau}: ví dụ cầm một cái cốc, người demo có lúc vòng tay qua bên trái vật cản, có lúc qua bên phải --- cả hai đều đúng. Một mạng hồi quy bị buộc tối thiểu hoá lỗi trung bình sẽ học ra \emph{giá trị trung bình} của "trái" và "phải", tức là đi \textbf{thẳng vào giữa} --- chính giữa vật cản. Đây gọi là vấn đề \textbf{multimodality} (đa chế độ), và nó là động lực trực tiếp khiến ngành chuyển từ "hồi quy ra 1 hành động" sang "\textbf{sinh} ra 1 hành động".

\subsubsection{Bước 2 --- "generative model" nghĩa là gì, nói không cần công thức}
Một generative model không học ra "một đáp án", nó học ra \textbf{cả một phân phối} các đáp án hợp lý, rồi mỗi lần dùng nó "rút mẫu" (sample) ngẫu nhiên một đáp án cụ thể từ phân phối đó. Ví dụ dễ hình dung: yêu cầu một mô hình sinh ảnh "vẽ một con mèo" 5 lần, mỗi lần ra 5 con mèo khác nhau --- mèo trắng, mèo đen, mèo nằm, mèo đứng --- không lần nào giống lần nào, nhưng lần nào cũng "đúng là mèo". Áp dụng cho robot: generative policy không trả về "hành động", nó trả về \textbf{một mẫu rút từ phân phối các hành động hợp lý} ứng với quan sát hiện tại --- lần này có thể vòng trái, lần sau (rút mẫu lại) có thể vòng phải, cả hai đều "hợp lý" theo dữ liệu huấn luyện.

\subsubsection{Bước 3 --- diffusion model sinh mẫu bằng cách nào}
Ý tưởng gốc (huấn luyện trên ảnh): lấy một bức ảnh thật, \textbf{làm nhiễu dần} (thêm một chút nhiễu Gauss) qua rất nhiều bước nhỏ, tới khi ảnh biến thành nhiễu trắng hoàn toàn --- gọi là \emph{forward process}. Mạng nơ-ron được huấn luyện để học \textbf{quá trình ngược}: xuất phát từ nhiễu trắng, \textbf{khử nhiễu dần} qua nhiều bước nhỏ để từng bước "làm lộ ra" một mẫu hợp lệ --- gọi là \emph{reverse/denoising process}. Lý do chia thành nhiều bước nhỏ: mỗi bước khử nhiễu là một bài toán \emph{dễ} (chỉ cần đoán ra một chút nhiễu vừa được thêm vào), nên ghép nhiều bước dễ lại, mạng học được cả một quá trình sinh mẫu rất phức tạp mà không cần học "nhảy thẳng" từ nhiễu ra kết quả cuối. \textbf{Diffusion Policy} chỉ là áp đúng cơ chế này lên hành động robot: denoiser nhận quan sát hiện tại (ảnh camera, vị trí khớp) làm điều kiện, và khử nhiễu ra không phải một ảnh, mà một \textbf{chuỗi hành động} tương lai.

\subsubsection{Bước 4 --- vì sao sinh nguyên "cả chuỗi" (chunk) một lần, không phải từng bước}
Nếu sinh từng hành động một, rồi mới sinh tiếp hành động kế (kiểu autoregressive), sai số nhỏ dễ tích luỹ và các hành động liên tiếp --- vì được sinh độc lập từng bước --- có thể "giật cục", không mượt theo thời gian. Sinh nguyên một \textbf{action chunk} (một chuỗi $H$ hành động) cùng lúc giải quyết cả hai vấn đề: hành động mượt/nhất quán hơn, và đỡ phải chạy lại toàn bộ mạng denoiser nhiều lần (tăng throughput). \textbf{Nhưng đây chính là cái giá phải trả mà \st{} phải xử lý}: một khi chunk đã được "chốt" (commit), nó là một gói quyết định cho cả $H$ bước tương lai cùng lúc --- không còn cơ hội sửa từng bước nhỏ như một bộ điều khiển cổ điển vẫn làm.

\subsubsection{Bước 5 --- flow matching khác diffusion ở đâu (ngắn gọn)}
Flow matching~\cite{lipman2023flow} theo đuổi cùng mục tiêu (sinh mẫu từ một phân phối phức tạp) nhưng thay vì "khử nhiễu ngẫu nhiên qua rất nhiều bước rời rạc", nó học một \textbf{luồng liên tục} (một trường vector, hình dung như một dòng chảy) di chuyển \emph{xác định} (một phương trình ODE, không ngẫu nhiên ở mỗi bước) từ nhiễu tới mẫu hợp lệ. Vì đường đi này mượt và xác định hơn, flow matching thường cần \emph{ít bước tích hợp hơn} để ra một mẫu so với diffusion. Nhưng về bản chất đối với \st{}, hai thứ giống nhau ở điểm quan trọng nhất: cả hai đều là một sampler nhiều tham số, chạy qua nhiều bước nội tại để ra một action chunk, và không ai viết được phương trình tường minh cho \emph{toàn bộ} quá trình đó.

\subsubsection{Bước 6 --- nối lại: vì sao chính những đặc điểm này khiến \st{} phải tồn tại}
Tổng hợp lại, diffusion/flow policy có ba đặc điểm khiến chúng \textbf{không thể verify hình thức theo cách cổ điển}: (a) là mạng nơ-ron cỡ chục--trăm triệu tham số; (b) phải chạy qua nhiều bước nội tại (10--100 bước khử nhiễu, hoặc nhiều bước tích hợp ODE) chỉ để ra \emph{một} action chunk; (c) có tính ngẫu nhiên/đa modal --- hỏi lại có thể ra chunk khác, đều "hợp lý" theo dữ liệu huấn luyện nhưng không ai đảm bảo hợp lý \emph{với tình huống an toàn cụ thể hiện tại}. Các công cụ neural-network reachability hiện có~\cite{ivanov2019verisig,tran2020nnv,huang2019reachnn} chỉ verify được mạng "khiêm tốn" qua horizon ngắn --- unroll một denoiser 100M+ tham số qua hàng chục bước, rồi qua cả contact dynamics vật lý, nằm ngoài khả năng đó rất xa. Đây chính là lý do \st{} chọn \textbf{không nhìn vào bên trong sampler}: coi $\pi_\theta$ là hộp đen chỉ đề xuất, rồi chỉ verify \emph{hậu quả} của một chunk cụ thể đã sinh ra --- đúng insight đã nêu ở \S\ref{sec:direction}.

\subsection{Vậy \st{} "kế thừa và tiếp tục phát triển" cái gì cụ thể?}
Nhìn theo bảng trên, có một điểm rất dễ bị bỏ qua nhưng lại là mạch chính: \textbf{công cụ toán của control cổ điển (reachable set, invariant set, receding-horizon re-check) không biến mất khi ta chuyển sang policy học được --- nó chỉ đổi vai trò.} Ở MPC (dòng 2), công cụ đó \emph{trực tiếp sinh ra} hành động. Ở \st{} (dòng cuối), công cụ đó \emph{không sinh hành động nữa} --- nó chỉ đứng bên cạnh, quan sát hành động mà một sampler mạnh hơn (nhưng không giải thích được) vừa đề xuất, và quyết định "cho qua / bác bỏ / sửa nhẹ". Đây chính là câu trả lời cho "kế thừa nền tảng control trong policy": \st{} không phát minh lại reachability hay MPC, nó \textbf{chuyển vai trò của chúng từ "người lái" sang "người giám sát"}, để công cụ cũ vẫn dùng được ngay cả khi "người lái chính" giờ là một mạng nơ-ron sinh xác suất mà không ai viết được phương trình cho nó.

\textbf{Ẩn dụ dễ hình dung}: hãy tưởng tượng một huấn luyện viên bay (flight instructor) ngồi cạnh một học viên lái máy bay. Học viên (generative policy) đề xuất một loạt động tác lái sắp thực hiện. Huấn luyện viên không tự lái máy bay và không dạy lại học viên ngay lúc đó --- huấn luyện viên chỉ nhìn trước vài giây (short-horizon), tính xem "nếu học viên thực hiện đúng động tác này, máy bay sẽ ở đâu, có va vào núi không" (reachtube), và chỉ can thiệp (giữ cần lái, hoặc gợi ý chỉnh nhẹ) khi thấy rõ nguy hiểm cụ thể (counterexample) --- chứ không kiểm tra toàn bộ "tư duy" của học viên. \st{} là phiên bản toán học chính xác của huấn luyện viên này, áp cho một "học viên" là generative robot policy.

\subsection{Vì sao đây là hướng đang phát triển, không phải một bài toán đã xong}
Ngành robot học hiện đang đi từ policy nhỏ (một tác vụ, một môi trường) sang \textbf{policy nền tảng} (foundation-model-style, VLA --- vision-language-action) huấn luyện trên dữ liệu khổng lồ, càng lúc càng giống generative policy về quy mô và độ khó phân tích. Khi các policy này được đưa vào robot thật gần con người, câu hỏi "ai đứng ra bảo đảm an toàn khi bản thân policy không thể verify được" trở thành câu hỏi trung tâm của cả ngành --- không riêng \st{}. Nói cách khác, \st{} là một điểm dữ liệu (data point) mới nhất, chứ chưa phải điểm cuối, trên đúng một trục phát triển mà paper MPC trước của bạn cũng nằm trên đó: \textbf{làm sao giữ được bảo đảm kiểu control cổ điển khi đối tượng cần điều khiển ngày càng "học được" và ngày càng khó viết phương trình tay}.

%=====================================================================
\section{Nền tảng (Foundations)}\label{sec:foundations}
%=====================================================================
\st{} không tự phát minh một lĩnh vực mới; nó \emph{ghép} năm khối kiến thức đã tồn tại độc lập, theo một cấu hình chưa ai dùng cho generative robot policy. Hiểu rõ mỗi khối --- và biên giới của nó --- là điều kiện để hiểu vì sao ShortStop hoạt động và vì sao nó khó làm đúng.

\subsection{Khối 1 --- Generative robot policy (diffusion / flow matching)}
Diffusion Policy~\cite{chi2023diffusion} và các biến thể planning-as-inference~\cite{janner2022planning,ajay2023conditional}, cùng flow matching~\cite{lipman2023flow}, là kiến trúc chuẩn hiện nay để học chính sách điều khiển robot bằng imitation learning. Đặc điểm quan trọng nhất cho ShortStop: chính sách không sinh ra \emph{một} hành động mà một \textbf{action chunk} $a_{t:t+H}$ --- một chuỗi hành động tương lai được commit cùng lúc, giúp tăng tính nhất quán thời gian và throughput~\cite{zhao2023aloha}. Hệ quả: bất kỳ cơ chế an toàn nào muốn can thiệp phải quyết định \emph{trước khi} cả chuỗi được thực thi, không phải theo từng bước một như control cổ điển.

Về mặt tính toán, sampler này là một mạng denoising (U-Net 1D-conv hoặc transformer) hàng chục--hàng trăm triệu tham số, được unroll qua nhiều bước khử nhiễu (10--100 bước DDIM). Reachability analysis cho mạng nơ-ron~\cite{ivanov2019verisig,tran2020nnv,huang2019reachnn,everett2021reachability} chỉ scale tới mạng "khiêm tốn" và horizon ngắn --- không tới cỡ sampler này khi unroll qua contact dynamics. \textbf{Đây chính là ràng buộc nền tảng khiến ShortStop phải đứng ngoài $\pi_\theta$ và không bao giờ reach xuyên qua nó.}

\subsection{Khối 2 --- Reachability analysis (set propagation)}
Các công cụ propagation tập trạng thái cho hệ phi tuyến/hybrid (CORA~\cite{althoff2015cora}, Flow*~\cite{chen2013flowstar}) tính một \emph{reachtube} --- dãy tập $\{\hat R_k\}_{k=0}^H$ chứa mọi trạng thái model-consistent. ShortStop dùng công thức đệ quy
\[
\hat R_{k+1} = \hat f(\hat R_k, a_{t+k}) \oplus B(\varepsilon_k + \bar w),
\]
biểu diễn bằng interval hoặc zonotope, với $\oplus$ là Minkowski sum. Điểm mấu chốt: ShortStop \textbf{không} reach qua $\pi_\theta$ (như neural-feedback-loop reachability~\cite{everett2021reachability} làm với một controller nơ-ron cố định) --- nó chỉ reach qua \emph{môi trường} phản ứng với một chunk \emph{đã được lấy mẫu cụ thể}. Đây là lý do bài toán trở nên tractable: một khi $a_{t:t+H}$ đã cố định, hệ không còn policy trong loop, chỉ còn một hệ động lực bị điều khiển (controlled system), và set propagation chuẩn áp dụng được.

\subsection{Khối 3 --- Signal Temporal Logic (STL) và robustness}
STL~\cite{maler2004monitoring} trang bị ngữ nghĩa định lượng (robustness)~\cite{donze2010robust,fainekos2009robustness} cho đặc tả thời gian: một giá trị $\rho$ dương nghĩa là thoả mãn với biên độ $\rho$. ShortStop mở rộng ý tưởng \emph{robust online monitoring}~\cite{deshmukh2017robust} --- thường được tính trên \emph{một} tín hiệu quan sát --- sang tính trên \emph{cả một reachtube}:
\[
\rho(\phi,\hat R) = \min_{0\le k\le H} \inf_{x\in\hat R_k} \mathrm{dist}(x,X_u).
\]
Vì $\hat R_k$ chứa mọi trạng thái khả dĩ, $\rho(\phi,\hat R)$ là một \textbf{lower bound chứng nhận} (certified lower bound) chứ không phải một điểm số nominal --- đây là khác biệt ngữ nghĩa quan trọng nhất so với STL-monitor cổ điển, và cũng là nơi literature 2022--2024 (model predictive monitoring, xem \S\ref{sec:related}) đã đi trước một phần.

\subsection{Khối 4 --- Counterexample-guided synthesis/falsification (CEGIS/CEGAR)}
CEGIS~\cite{solar2006combinatorial} và CEGAR~\cite{clarke2000counterexample} là nguyên lý tổng quát: dùng một \emph{witness cụ thể} (không chỉ một điểm số) để (a) chứng minh sự bác bỏ và (b) định hướng vòng sửa kế tiếp. Các công cụ falsification cho hệ cyber-physical (S-TaLiRo~\cite{annpureddy2011staliro}, Breach~\cite{donze2010breach}, VerifAI~\cite{dreossi2019verifai}) áp dụng nguyên lý này để tìm input làm rớt một đặc tả. ShortStop \emph{nhập khẩu} nguyên lý CEGIS vào bài toán lọc chunk tại runtime: bài toán đối nghịch
\[
(k^\star, x^\star) = \arg\min_{0\le k\le H,\, x\in\hat R_k} \mathrm{dist}(x,X_u)
\]
giải bằng projected gradient / vertex sampling trên zonotope, và $x^\star$ vừa là chứng cứ bác bỏ vừa cho gradient sửa $d$.

\subsection{Khối 5 --- Predictive safety filters, MPC, CBF (nền tảng gần nhất với kinh nghiệm cá nhân)}
Predictive safety filter~\cite{wabersich2021predictive}, model-predictive shielding~\cite{bastani2021safe}, và control barrier function QP~\cite{ames2017cbf,ames2019cbf-ecc,cheng2019end} là các baseline gần nhất với nền tảng MPC/QP đã có sẵn (từ paper MPC trước của tác giả). Điểm khác biệt cấu trúc quan trọng: các phương pháp này (i) thường sửa \emph{một} input tại một bước, và (ii) cần một mô hình optimal-control tường minh hoặc một barrier thiết kế tay. ShortStop chọn/sửa cả một \emph{batch} $K$ chunk sinh từ sampler, và chứng nhận một đặc tả STL đa dạng hơn chỉ set-invariance.

\medskip
\noindent\textbf{Tổng kết Sec.~\ref{sec:foundations}:} Khối 1 và 5 là nơi bài báo \emph{dùng} kiến thức nền có sẵn (generative policy như một black-box, MPC/CBF như baseline). Khối 2--4 (reachability, STL-over-tube, CEGIS) là nơi bài báo phải \emph{lắp ráp lại} theo một cấu hình mới --- và đây cũng chính xác là ba module nặng nhất trong \emph{ShortStop\_Research\_Plan.tex} (Module 3.D--F).

%=====================================================================
\section{Hướng hoạt động: insight cốt lõi}\label{sec:direction}
%=====================================================================
\subsection{Bài toán bất khả thi $\to$ bài toán khả thi}
Câu mở đầu Sec.~I của bài báo nêu chính xác động lực: \emph{"Full formal verification of such policies is out of reach [\ldots] Yet a crucial observation drives this paper: we do not need to verify the policy to make it safe. We only need to verify, cheaply and just-in-time, the consequences of the chunk it is about to execute, over the short horizon during which that chunk is authoritative."}

Đây là một phép \textbf{tách bài toán} (problem decoupling) rất đặc trưng của formal-methods thực dụng: thay vì hỏi "$\pi_\theta$ có luôn an toàn với mọi observation không?" (một câu hỏi $\forall$-lượng từ hoá trên không gian tham số 100M+ chiều, bất khả thi), ShortStop hỏi "\emph{với chunk cụ thể mà $\pi_\theta$ vừa đề xuất}, hậu quả ngắn hạn của nó có an toàn không?" (một câu hỏi trên một hệ động lực đã cố định input, khả thi bằng set propagation cổ điển). Về bản chất, $\pi_\theta$ bị đối xử như một \textbf{black-box proposer} chứ không phải đối tượng cần verify.

\subsection{Vì sao "ngắn hạn" (short-horizon) là chìa khoá, không phải một sự đánh đổi tầm thường}
Ba lý do kỹ thuật khiến short-horizon vừa \emph{đủ} vừa \emph{cần}:
\begin{itemize}
  \item \textbf{Đủ}: vì thực thi la receding-horizon --- chỉ hành động đầu của chunk được commit, certification được làm mới lại mỗi bước (Sec.~III.G) --- nên "an toàn trong $H$ bước tới, chứng nhận lại mỗi bước" tương đương an toàn liên tục, miễn Corollary~1 (recurrent feasibility) giữ được.
  \item \textbf{Cần} (theo Prop.~1): tập reachable nở ra theo cấp $L^k$ (Lipschitz) mỗi bước --- khiến horizon dài làm nổ conservatism (tỉ lệ reject) trước khi mất tính đúng đắn. Do đó tồn tại một $H^\star$ tối ưu (thực nghiệm $H^\star=8$), không phải "dài hơn luôn tốt hơn".
  \item \textbf{Không cần reach qua $\pi_\theta$}: vì chunk đã lấy mẫu cụ thể, chi phí mỗi bước độc lập với kích thước mạng $\pi_\theta$ --- đây là lý do ShortStop scale tới sampler 100M+ tham số mà các công cụ neural-reachability không scale tới được.
\end{itemize}

\subsection{Bốn giai đoạn P-R-C-S như một vòng lặp defense-in-depth}
\begin{description}
  \item[(P) Propose] Lấy $K=8$ chunk mẫu $\{a^{(i)}\}\sim\pi_\theta(\cdot\mid o_t)$ --- không tìm kiếm toàn không gian hành động, chỉ \emph{lọc} trong các mẫu policy tự đề xuất.
  \item[(R) Reach] Lan truyền reachtube sound cho từng chunk (Eq.~1).
  \item[(C) Certify] Tính STL robustness-to-go lower bound $\rho(\phi,\hat R)$ (Eq.~2) trên cả tube.
  \item[(S) Select/Repair] Tìm counterexample (Eq.~3); nếu có, sửa theo hướng gradient $d$ (Eq.~4) trong trust region $\delta$; chọn chunk admissible tối đa task-progress $g$ (Eq.~5); nếu tập admissible trống, gọi fallback $\pi_{fb}$.
\end{description}
Cấu trúc này là một dạng \textbf{defense-in-depth}: (R)+(C) là lớp phòng thủ "rẻ nhưng có thể quá conservative" (bao trọn cả tube); (S) là lớp phòng thủ "đắt hơn nhưng chính xác" (chỉ can thiệp khi có witness cụ thể) --- ablation Table~III cho thấy chính lớp CEGIS này đóng góp lớn nhất (violation $2.9\%\to1.7\%$) vì nó "bắt" được các vi phạm ở góc tube mà interval bound làm tròn mất.

%=====================================================================
\section{Cách hiện thực: toán học, thuật toán, và bộ ba định lý}\label{sec:implementation}
%=====================================================================
\subsection{Mô hình bài toán}
Trạng thái thật $x_{t+1}=f(x_t,a_t)+w_t$, $\|w_t\|\le\bar w$; $\pi_\theta$ sinh $a_{t:t+H}\sim\pi_\theta(\cdot\mid o_t)$; đặc tả an toàn là STL $\phi$ (điển hình: avoidance $\square_{[0,H]}(x\notin X_u)$, có thể kết hợp progress $\lozenge_{[0,H]}$).

\subsection{Ba giả thiết soundness --- nơi "certificate" thật sự bắt nguồn}
\begin{itemize}
  \item \textbf{Giả thiết 1} (reachable-set sound): $x_t\in\hat R_0$ và $\|f-\hat f\|\le\varepsilon(x,a)$ $\Rightarrow$ $x_{t+k}\in\hat R_k$ với mọi $k\le H$. Đây là hợp đồng soundness tiêu chuẩn của mọi công cụ reachability~\cite{althoff2015cora,chen2013flowstar}; $\varepsilon$ và $\bar w$ là "cái giá" phải trả để giữ hợp đồng này đúng.
  \item \textbf{Giả thiết 2} (STL bound sound): $\rho(\phi,\hat R)\le\rho(\phi,\xi,0)$ cho mọi $\xi$ nằm trong tube --- đúng theo cấu trúc với ngữ nghĩa interval STL~\cite{deshmukh2017robust}.
  \item \textbf{Giả thiết 3} (fallback tồn tại): có $\pi_{fb}$ và tập robust-control-invariant $X_s\subseteq\{x:\mathrm{dist}(x,X_u)\ge\epsilon\}$ sao cho từ $X_s$, $\pi_{fb}$ giữ trạng thái trong $X_s$ dưới mọi nhiễu $\|w\|\le\bar w$ --- tiền đề chuẩn của predictive shielding~\cite{wabersich2021predictive,bastani2021safe}.
\end{itemize}

\subsection{Ba định lý/mệnh đề, và ý tưởng chứng minh}
\begin{theorem}[An toàn hữu hạn thời gian --- Thm.~1]
Dưới Giả thiết 1--2, nếu ShortStop commit chunk $a^\star$ với $\rho(\phi,\hat R^\star)\ge\epsilon\ge 0$, thì quỹ đạo closed-loop thực tế thoả $\phi$ trên $[t,t+H]$, và duy trì biên an toàn $\mathrm{dist}(x_{t+k},X_u)\ge\epsilon$.
\end{theorem}
\emph{Ý tưởng}: Giả thiết 1 $\Rightarrow$ quỹ đạo thật nằm trong tube; Giả thiết 2 $\Rightarrow$ bound chứng nhận là lower bound của robustness thật của quỹ đạo đó; bound $\ge\epsilon\ge 0$ $\Rightarrow$ thoả mãn với biên $\epsilon$. Đây thuần là một \textbf{chuỗi bất đẳng thức lồng nhau} ($\rho(\phi,\xi,0)\ge\rho(\phi,\hat R^\star)\ge\epsilon$) --- đơn giản về kỹ thuật chứng minh, nhưng \emph{sức mạnh} của nó nằm ở việc hai giả thiết nền (1 và 2) đã gộp toàn bộ độ khó formal-verification vào một cam kết calibration ($\varepsilon,\bar w$) có thể kiểm định bằng dữ liệu, thay vì phải verify $\pi_\theta$.

\begin{corollary}[An toàn mọi thời điểm --- Cor.~1]
Thêm Giả thiết 3, nếu hệ khởi đầu trong $X_s$ và tại mọi bước hoặc (a) commit chunk chứng nhận biên $\ge\epsilon$, hoặc (b) không có chunk admissible và dùng $\pi_{fb}$, thì $x_t\in X_s$ với mọi $t\ge 0$.
\end{corollary}
\emph{Ý tưởng}: quy nạp theo $t$ trên bất biến $x_t\in X_s$; case (a) dùng Thm.~1 để giữ margin $\ge\epsilon$ (nên successor vẫn trong $X_s$); case (b) dùng trực tiếp Giả thiết 3. Đây là kỹ thuật \textbf{switched-system invariance} tiêu chuẩn, nhưng lần đầu áp dụng cho một switching logic dựa trên certificate reachtube+STL+CEGIS (chứ không phải dựa trên một QP hay automaton).

\begin{proposition}[Trade-off conservatism--horizon --- Prop.~1]
Với $L$ là hằng số Lipschitz cục bộ của $\hat f$, bán kính tube $r_k$ thoả (Grönwall rời rạc)
\[
r_k \le L^k r_0 + (\varepsilon_{\max}+\bar w)\frac{L^k-1}{L-1}\quad (L>1),
\]
và tỉ lệ reject $\alpha(H)=\Pr[\rho(\phi,\hat R)<\epsilon]$ thoả
\[
\alpha(H) \le \alpha^\star + \beta(\epsilon+r_H),
\]
với $\alpha^\star=\Pr[\rho_{nom}<0]$ là tỉ lệ chunk \emph{thực sự} không an toàn, và $\beta$ chặn mật độ robustness gần ngưỡng.
\end{proposition}
\emph{Ý tưởng}: tách sự kiện reject thành hai phần --- (i) chunk thực sự nguy hiểm ($\alpha^\star$, không thể tránh) và (ii) chunk nằm trong "băng bảo thủ" $[0,\epsilon+r_H)$ tạo ra bởi margin và độ nở của tube; chặn mật độ trong băng đó bằng $\beta$. Đây là đóng góp \textbf{toán học mới nhất} của bài báo: nó biến trực giác định tính "horizon dài hơn = certify nhiều hơn nhưng reject nhiều hơn" thành một bất đẳng thức định lượng, và \emph{tiên đoán được} sự tồn tại của $H^\star$ nội tại --- điều được xác nhận thực nghiệm ở Fig.~4 ($H^\star=8$).

\subsection{Thuật toán và độ phức tạp}
Algorithm~1 (trong bài báo) chạy $K$ nhánh độc lập (mỗi nhánh: 1 lần Reach + $M$ vòng CE-search), sau đó chọn chunk admissible tối đa $g(\cdot)$. Với $K=8,H=8,M=16$: latency trung vị $8.9$~ms (Reach $3.1$ms + STL-to-go $1.4$ms + CE-search $3.0$ms + Repair $1.4$ms), nằm trong budget $10$--$20$Hz. Điểm hiện thực quan trọng: toàn bộ $K$ nhánh \emph{song song hoá được} và chi phí \textbf{không phụ thuộc kích thước $\pi_\theta$} --- đây là lý do phương pháp scale.

%=====================================================================
\section{Cái dễ và cái khó}\label{sec:hard-easy}
%=====================================================================
\subsection{Cái dễ (tương đối) --- những phần dựa trên formal-methods đã trưởng thành}
\begin{itemize}
  \item \textbf{Set propagation trên interval/zonotope 2D--6D}: công thức (1) là chuẩn CORA/Flow*~\cite{althoff2015cora,chen2013flowstar}; tự code interval arithmetic cho 2D (như \emph{Research\_Plan} Phase~1 đề xuất) là bài tập kỹ thuật, không phải nghiên cứu mới.
  \item \textbf{STL robustness đệ quy trên interval}: công thức (2) và các quy tắc $\wedge,\vee,\square,\lozenge$ trên bound là recursion tiêu chuẩn~\cite{deshmukh2017robust,fainekos2009robustness}.
  \item \textbf{MPC-Filter và CBF-Shield baseline}: tận dụng trực tiếp kinh nghiệm QP (Hildreth) đã có sẵn từ paper MPC trước --- đây chính xác là lý do \emph{ShortStop\_Research\_Plan.tex} đánh dấu Sec.~II.D là \bridge{} (cầu nối) chứ không phải \new{}.
  \item \textbf{Receding-horizon re-certification}: về mặt vòng lặp điều khiển, giống MPC cổ điển --- không có gì mới trong \emph{cấu trúc lặp lại mỗi bước}.
  \item \textbf{Ablation/thống kê thực nghiệm} (bootstrap, seeds, metrics): công việc kỹ thuật thực nghiệm chuẩn, không đòi hỏi lý thuyết mới.
\end{itemize}

\subsection{Cái khó --- những phần đòi hỏi đóng góp/khéo léo thật sự}
\begin{itemize}
  \item \textbf{Calibrate $\varepsilon$ đúng cách}: $\varepsilon$ là "linh hồn" của Giả thiết 1 --- nếu high-quantile residual bound bị đánh giá thấp (đặc biệt dưới distribution shift, ví dụ novel contact mode), \emph{toàn bộ} certificate sập mà không có cảnh báo hình thức nào (bài báo tự nhận đây là failure mode chính, Sec.~VIII.b). Đây là nơi ranh giới giữa "chứng minh được" và "empirically sound" (Fig.~5) trở nên mong manh.
  \item \textbf{Xây $X_s$ và $\pi_{fb}$ hợp lệ (Giả thiết 3)}: không tự động đúng --- phải thiết kế một braking/retreat controller \emph{và} một robust-control-invariant set \emph{tương thích} với nó cho mỗi môi trường/robot. Đây là công việc control-theory truyền thống nhưng tốn công stitching vào pipeline generative.
  \item \textbf{Giải bài toán đối nghịch (3) đủ nhanh và đủ tốt}: projected gradient/vertex sampling trên zonotope phải vừa rẻ (để giữ latency 8.9ms) vừa đủ mạnh để không bỏ lọt counterexample (nếu bỏ lọt, certificate optimistic một cách âm thầm). Đây là một bài toán tối ưu non-convex nhỏ nhưng "phải đúng, không chỉ nhanh".
  \item \textbf{Chứng minh Grönwall-type cho Prop.~1}: đòi hỏi kỹ năng chứng minh recursion tuyến tính + tách xác suất theo băng conservatism --- loại toán không có trong nền tảng MPC/QP trước đó của tác giả (được \emph{Research\_Plan} đánh dấu \new{} và dành riêng 2--3 buổi).
  \item \textbf{Cân bằng safety vs. liveness}: Corollary~1 chỉ đảm bảo an toàn, không đảm bảo tiến triển --- 22\% tình huống reject không recover được (rơi vào fallback). Thiết kế repair loop (Eq.~4) để giảm tỉ lệ này (recovery $0.44\to0.78$ sau khi thêm Repair) là phần \emph{engineering-cho-đúng-mục-tiêu} khó nhất về mặt thực nghiệm.
  \item \textbf{Tính đúng đắn dưới black-box, high-dim, multimodal sampler}: khác các baseline CBF/MPC (giả định model tường minh, single input), ShortStop phải certify một \emph{batch} chunk từ một sampler mà chính nó không được reach qua --- buộc phải tách rời hoàn toàn "chứng minh" (trên môi trường) khỏi "sinh đề xuất" (trên policy), một sự phân tách concept không hiển nhiên.
  \item \textbf{Đánh giá trên 4 testbed + 3 stress-test family + 6 baseline + 7 metric}: khối lượng thực nghiệm lớn, đặc biệt phần LIBERO/ManiSkill2/RLBench hiện tại là \emph{model-derived projections} (không phải full closed-loop run) --- bài báo tự nhận đây là giới hạn (Sec.~V.F, VIII.f) và đây chính là phần "khó nhất về resource" nếu ai muốn tái lập đầy đủ (cần train diffusion policy thật + GPU + simulator nặng).
\end{itemize}

%=====================================================================
\section{Vì sao đây là một đóng góp mới và khó}\label{sec:novelty}
%=====================================================================
\subsection{Mới ở đâu — không phải một thành phần, mà là một tổ hợp chưa ai làm}
Bài báo tự định vị rất rõ ở Sec.~II.H (Positioning/gap): máy an toàn hiện có chỉ làm được \emph{một trong bốn} điều, không bao giờ cả bốn:
\begin{enumerate}[label=(\roman*)]
  \item sửa/retrain generator (SafeDiffuser~\cite{xiao2023safediffuser}, sampler-projection~\cite{power2023sampling}) --- ShortStop \emph{không} retrain;
  \item chỉ chứng nhận \emph{một} action đã sửa, không phải một \emph{batch} chunk temporal-extended (MPC-Filter~\cite{wabersich2021predictive}, CBF-QP~\cite{ames2017cbf}) --- ShortStop chứng nhận cả batch;
  \item cần một CBF thiết kế tay hoặc một model optimal-control tường minh và tractable --- ShortStop chỉ cần $\hat f,\varepsilon$ (mô hình học được hoặc tuyến tính hoá cục bộ);
  \item chỉ score một quỹ đạo nominal, không có tube certificate hay counterexample (STL-monitor thuần~\cite{deshmukh2017robust}) --- ShortStop certify cả tube \emph{và} chủ động tìm counterexample.
\end{enumerate}
Cái mới không nằm ở bất kỳ khối riêng lẻ (reachability, STL, CEGIS đều là kỹ thuật cũ) mà ở việc \textbf{ba khối này lần đầu được ghép thành một vòng lặp per-decision, áp dụng cho một generative sampler tần suất cao, với một chứng minh toán học đi kèm} --- và counterexample "làm việc gấp đôi" (vừa là chứng cứ bác bỏ, vừa là hướng sửa) là điểm khác biệt tinh tế nhất so với việc chỉ dùng reachability+STL như một bộ lọc nhị phân.

\subsection{Khó ở đâu — ba loại khó khác nhau, không chỉ một}
\begin{itemize}
  \item \textbf{Khó về mặt khái niệm (conceptual)}: phải tìm đúng "lát cắt" của bài toán để tách phần \emph{intractable} (verify $\pi_\theta$) khỏi phần \emph{tractable} (verify hậu quả của 1 chunk cố định) --- một sai lầm phổ biến là thử certify cả chuỗi denoising, dẫn ngay về bài toán bất khả thi ban đầu.
  \item \textbf{Khó về mặt toán học (formal)}: chứng minh phải giữ được tính đúng đắn dưới \emph{ba} nguồn bất định cùng lúc --- model error $\varepsilon$, disturbance $\bar w$, và tính ngẫu nhiên của sampler $\pi_\theta$ --- mà không cần biết gì về cấu trúc nội tại của $\pi_\theta$. Đây là lý do proof phải đặt toàn bộ gánh nặng vào 3 giả thiết soundness \emph{bên ngoài} $\pi_\theta$, và Prop.~1 phải dùng Grönwall để lượng hoá trade-off một cách tường minh (không chỉ định tính).
  \item \textbf{Khó về mặt hệ thống/kỹ thuật (systems)}: toàn bộ vòng lặp (K reachtube + K.M inner-search) phải chạy trong <10ms cạnh một GPU policy, và phải \emph{đúng} (không optimistic) dù chạy nhanh --- một sự đánh đổi kinh điển giữa soundness và tốc độ mà bài báo giải quyết bằng cách giữ set propagation ở không gian \emph{low-dimensional} (toạ độ an toàn task-space, không phải action-space đầy đủ hay latent space của $\pi_\theta$).
\end{itemize}

\subsection{Giới hạn của tính "mới" — cần nhìn thẳng}
Một phần "mới" của bài báo (STL robustness tính trên reachtube/predicted set thay vì một quỹ đạo nominal) \emph{đã có tiền lệ gần} trong literature model-predictive-monitoring (xem \S\ref{sec:related}, bài Yu et al.) công bố 2022--2024, dùng đúng ý tưởng "dùng model để dự đoán reachable set rồi evaluate STL trên đó, thay vì trên 1 trajectory". Sự khác biệt thực sự của ShortStop với dòng này là: (a) áp dụng cho \emph{lựa chọn hành động của một generative sampler} chứ không phải \emph{giám sát bị động} một hệ đã biết luật điều khiển; và (b) đóng thêm một vòng lặp CEGIS chủ động (reject-and-\emph{repair}), không chỉ certify/falsify. Nhìn nhận đúng mức độ mới này (kết hợp + áp dụng mới, không phải phát minh nguyên lý STL-trên-reachable-set) là quan trọng để đánh giá công bằng.

%=====================================================================
\section{Bốn bài báo liên quan --- điểm mạnh/yếu so với ShortStop}\label{sec:related}
%=====================================================================
Bốn bài được chọn theo đúng yêu cầu: bắt đầu từ reference gốc của ShortStop (SafeDiffuser [24], MPC predictive safety filter [33], STL robust online monitoring [53]), sau đó tìm rộng ra literature 2022--2025 (model predictive monitoring cho STL; và một bài đối thủ trực tiếp công bố sau ShortStop's reference list, 2025) để xem ai đang cạnh tranh/tiếp nối ý tưởng gần nhất.

\subsection{Bài 1 --- SafeDiffuser (Xiao, Wang, Gan, Rus)}
\textbf{Nguồn}: \emph{"SafeDiffuser: Safe Planning with Diffusion Probabilistic Models"}, arXiv:2306.00148 (2023), ICLR 2025~\cite{xiao2023safediffuser}. Đây là reference [24] trong ShortStop.

\textbf{Phương pháp}: nhúng control barrier functions trực tiếp vào quá trình denoising của diffusion model, tạo ra "finite-time diffusion invariance" --- nói cách khác, \emph{sửa chính sampler} để các bước khử nhiễu trung gian bị ràng buộc ở lại trong tập an toàn, thay vì lọc output sau khi sinh xong. Thử nghiệm trên maze path generation, legged locomotion, và 3D manipulation.

\textbf{Điểm mạnh so với ShortStop}: (a) không cần vòng lọc runtime riêng biệt --- an toàn được "dệt" vào chính quá trình sinh, nên không có overhead tách rời kiểu Reach+CE-search; (b) không cần một mô hình động lực $\hat f,\varepsilon$ tường minh ở runtime vì ràng buộc được áp trong không gian denoising.

\textbf{Điểm yếu so với ShortStop}: (a) phải \emph{sửa/huấn luyện lại} sampler --- không áp dụng được cho một $\pi_\theta$ đã đóng băng (frozen, đã deploy) như ShortStop yêu cầu; (b) guarantee gắn với xấp xỉ của quá trình diffusion invariance, không phải một per-decision certificate độc lập với việc training có đúng không; (c) không xử lý một batch $K$ chunk với lựa chọn/so sánh task-progress --- an toàn và tính task-optimal bị trộn vào cùng một quá trình sinh, khó tách bạch "cái gì làm rớt task success vì lý do an toàn"; (d) không có bằng chứng loại bounded-time/all-time theorem tách biệt theo kiểu ShortStop, cũng không có counterexample-guided repair tường minh.

\subsection{Bài 2 --- Predictive Safety Filter (Wabersich \& Zeilinger)}
\textbf{Nguồn}: \emph{"A predictive safety filter for learning-based control of constrained nonlinear dynamical systems"}, Automatica 129:109597 (2021)~\cite{wabersich2021predictive}. Reference [33] --- chính là baseline \texttt{MPC-Filter} trong Table~II của ShortStop.

\textbf{Phương pháp}: biến một hệ có ràng buộc thành một hệ "không ràng buộc an toàn" bằng một bộ lọc MPC — nhận input đề xuất từ policy RL, và tối thiểu-hoá sửa nó qua một constrained optimal-control problem dựa trên data-driven model, cập nhật liên tục một "safety policy".

\textbf{Điểm mạnh so với ShortStop}: (a) lý thuyết recursive-feasibility rất trưởng thành (kế thừa MPC cổ điển~\cite{mayne2000constrained}); (b) chỉ cần sửa tối thiểu 1 action, chi phí tính toán/độ phức tạp triển khai thường thấp hơn khi model tuyến tính/lồi; (c) không cần một sampler đề xuất nhiều candidate — hoạt động tốt với policy sinh 1 action/bước.

\textbf{Điểm yếu so với ShortStop} (theo đúng số liệu Table~II): violation cao hơn ShortStop $4$--$5\times$ ($4.9\%$ vs $1.2\%$), success thấp hơn ($65.7\%$ vs $68.4\%$), precision thấp hơn nhiều ($0.64$ vs $0.94$), và latency \emph{chậm hơn} ($12.4$ms vs $8.9$ms) — vì QP phải giải cho mỗi input thay vì tái sử dụng gradient counterexample. Về cấu trúc: (a) chỉ tác động lên \emph{một} action, không có khái niệm batch chunk temporal-extended; (b) cần một model optimal-control tractable (thường tuyến tính/lồi hoá) — kém phù hợp khi hành vi generative đa modal (nhiều chế độ hành động khả dĩ mà một QP đơn không nắm bắt); (c) không có STL hay đặc tả thời gian phong phú, chỉ constraint set-invariance dạng chuẩn.

\subsection{Bài 3 --- Model Predictive Monitoring of STL (Yu, Dong, Yin, Li)}
\textbf{Nguồn}: \emph{"Model Predictive Monitoring of Dynamical Systems for Signal Temporal Logic Specifications"}, arXiv:2209.12493, Automatica (2024)~\cite{yu2024modelpredictive}. Đây \emph{không} nằm trong reference gốc của ShortStop nhưng là "tổ tiên" gần nhất về mặt ý tưởng của Sec.~III.E (STL robustness-to-go trên reachtube) — tìm được bằng cách mở rộng từ nhánh II.F (robust online STL monitoring [53]) sang literature 2022--2024.

\textbf{Phương pháp}: giám sát online một STL formula bằng cách, tại mỗi bước, dùng mô hình động lực \emph{đã biết} để tính \emph{feasible set} (tập trạng thái đạt được dưới control input chưa biết), rồi evaluate STL trên tập đó — cho phép "falsify hoặc certify đặc tả \emph{trước} so với thuật toán không dùng model", chứ không chỉ đánh giá trên quỹ đạo quan sát được tới hiện tại.

\textbf{Điểm mạnh so với ShortStop}: (a) đây chính xác là ý tưởng "STL trên reachable set thay vì trên 1 trajectory" mà ShortStop dùng ở Eq.~2 — nghĩa là bài này đã thiết lập \emph{lý thuyết giám sát} nền tảng trước, với phân tích chặt hơn về early-detection; (b) là một bài monitoring "thuần", tổng quát hơn cho nhiều hệ cyber-physical, không giới hạn ở robot manipulation.

\textbf{Điểm yếu so với ShortStop}: (a) đây là một \textbf{monitor bị động} — phát hiện vi phạm sớm, nhưng không đóng vòng lặp \emph{hành động}: không có khái niệm reject-and-\emph{repair}, không chọn giữa nhiều candidate, không tích hợp vào một policy loop; (b) không nhằm vào generative/stochastic sampler cao chiều — giả định luật điều khiển \emph{chưa biết nhưng đơn nhất} tại mỗi bước, không phải một phân phối đa modal trên hành động; (c) không có counterexample-guided search kiểu CEGIS để vừa chứng minh reject vừa cho hướng sửa — chỉ dừng ở "chứng nhận/bác bỏ". Nói cách khác: bài 3 cung cấp \emph{nửa đầu} (Reach + Certify) của vòng lặp P-R-C-S, và ShortStop identify đúng phần còn thiếu (Select/Repair chủ động bằng CEGIS) là chỗ tạo ra phần lớn lợi ích thực nghiệm (ablation Table~III: CE-search + Repair đóng góp lớn nhất).

\subsection{Bài 4 --- Path-Consistent Safety Filtering for Diffusion Policies (2025, đối thủ trực tiếp/đồng thời)}
\textbf{Nguồn}: \emph{"From Demonstrations to Safe Deployment: Path-Consistent Safety Filtering for Diffusion Policies"}, arXiv:2511.06385 (2025)~\cite{pathconsistent2025}. Tìm được khi mở rộng tra cứu ra khỏi reference gốc, tìm bài \emph{gần nhất về mục tiêu} (an toàn cho diffusion policy tại runtime) công bố sau ShortStop.

\textbf{Phương pháp}: một bộ lọc an toàn hậu kiểm (post-hoc) cho diffusion policy, giữ "path-consistency" — nghĩa là chunk sau khi lọc vẫn nhất quán với phân phối mà diffusion model đã học, tránh việc sửa quá tay khiến hành động rời khỏi phân phối huấn luyện. Hoạt động như một wrapper tương thích với kiến trúc diffusion-policy bất kỳ.

\textbf{Điểm mạnh so với ShortStop}: (a) mục tiêu "path-consistency" giải quyết trực diện một rủi ro mà ShortStop chỉ giảm nhẹ gián tiếp qua trust region $\delta$ ở Eq.~4 — nếu bước sửa của ShortStop đẩy chunk quá xa khỏi phân phối $\pi_\theta$, hành vi có thể trở nên bất thường dù vẫn "an toàn" theo nghĩa hình học; bài 4 đặt việc giữ tính hợp lý theo phân phối lên làm mục tiêu chính, không phải ràng buộc phụ; (b) cùng nhắm đúng vào diffusion policy (không tổng quát hoá cho flow-matching như ShortStop), nên có thể tối ưu sâu hơn cho đặc thù kiến trúc denoiser.

\textbf{Điểm yếu so với ShortStop} (dựa trên các mô tả công khai — cần đọc full-text để xác nhận chi tiết định lượng): (a) không thấy công bố một bộ ba định lý dạng bounded-time / all-time / conservatism-horizon trade-off tương ứng Thm.~1/Cor.~1/Prop.~1 của ShortStop — nghĩa là guarantee có vẻ dừng ở mức "constraint satisfaction" hơn là một chuỗi chứng minh lượng hoá được sự đánh đổi theo horizon; (b) không thấy cơ chế counterexample-guided search tường minh (tìm witness cụ thể trong reachable set để vừa chứng minh reject vừa cho hướng sửa) — có vẻ nghiêng về ràng buộc/tối ưu trực tiếp trên trajectory hơn là một vòng CEGIS; (c) chỉ nhắm diffusion policy, còn ShortStop tuyên bố tổng quát cho \emph{bất kỳ} generative chunk-proposer kể cả flow-matching (đã thử ablation). \emph{Ghi chú khách quan}: đây là công trình \textbf{đồng thời/gần đồng thời} với ShortStop (cùng năm 2025, ShortStop nộp ICRA 2026) — nên xem là bằng chứng cho thấy "an toàn hậu kiểm cho diffusion policy" đang là một hướng nóng, chứ không dùng để khẳng định chắc chắn ShortStop "hơn" về mọi mặt mà không đọc kỹ full-text của bài 4.

\subsection{Bảng tổng hợp so sánh}
\begin{table}[H]
\centering\small
\begin{tabularx}{\textwidth}{@{}l X X X X@{}}
\toprule
Tiêu chí & SafeDiffuser~\cite{xiao2023safediffuser} & MPC-Filter~\cite{wabersich2021predictive} & MPM-STL~\cite{yu2024modelpredictive} & Path-Consistent~\cite{pathconsistent2025} \\
\midrule
Sửa generator? & Có (nhúng CBF vào denoiser) & Không & N/A (không có generator) & Không (post-hoc) \\
Đối tượng bảo vệ & 1 trajectory sinh ra & 1 action & 1 tín hiệu quan sát & 1 chunk diffusion \\
Dùng reachability? & Không (invariance trong denoising) & Có (rollout dự đoán) & Có (feasible/reachable set) & Không rõ (path-consistency, không phải tube) \\
STL / temporal spec? & Không & Không & Có (STL đầy đủ) & Không rõ \\
Counterexample-guided? & Không & Không & Không (monitor thuần) & Không rõ \\
Guarantee chính & Diffusion invariance (soft) & Recursive feasibility & Sound/early falsify-certify & Constraint satisfaction \\
Batch $K$ chunk + chọn? & Không & Không (1 input) & Không & Không rõ \\
\bottomrule
\end{tabularx}
\caption{So sánh 4 bài liên quan với cấu trúc P-R-C-S + 3-định-lý của ShortStop. "Không rõ" = thông tin không xác nhận được từ tóm tắt công khai, cần đọc full-text.}
\end{table}

%=====================================================================
\section{Tổng kết vị trí của ShortStop và câu hỏi còn mở}\label{sec:conclusion}
%=====================================================================
\subsection{Vị trí trong bức tranh literature}
Đặt cạnh 4 bài trên, ShortStop chiếm một điểm khá đặc thù: nó là bài \emph{duy nhất} trong nhóm này (a) \textbf{không sửa generator} (như MPC-Filter, MPM-STL), (b) \textbf{có đặc tả STL phong phú trên cả tube} (như MPM-STL, khác SafeDiffuser/MPC-Filter/Path-Consistent), (c) \textbf{đóng vòng CEGIS chủ động} (khác cả 4 bài), và (d) \textbf{chứng nhận một batch chunk rồi chọn/sửa} (khác MPC-Filter/CBF vốn chỉ sửa 1 input). Nói cách khác, ShortStop nằm ở giao của ba trục: \emph{post-hoc} (như MPC-Filter, Path-Consistent) $\times$ \emph{reachability+STL rigor} (như MPM-STL) $\times$ \emph{counterexample-guided} (chưa ai trong 4 bài có) — và chính giao điểm ba trục này, chứ không phải một trục riêng lẻ, là nơi đóng góp thật của bài báo.

\subsection{Những câu hỏi còn mở (đáng đọc tiếp / có thể là hướng mở rộng)}
\begin{itemize}
  \item Guarantee của ShortStop là \emph{conditional} trên $\varepsilon$ đúng — chuyển sang một bound xác suất (conformal calibration) sẽ nối trực tiếp với literature "robust conformal prediction cho STL runtime verification dưới distribution shift" (CPS 2024) mà tra cứu mở rộng đã chạm tới.
  \item So sánh trực tiếp, định lượng với bài Path-Consistent (2511.06385) trên cùng benchmark là việc chưa làm (và không thể làm bằng suy luận từ abstract) — cần đọc full-text/hoặc chạy lại cả hai trên cùng LIBERO/ManiSkill2 để có kết luận công bằng.
  \item Câu hỏi liveness (22\% reject không recover) chưa có lý thuyết đi kèm — đây là khoảng trống lý thuyết còn lại rõ nhất, khác với an toàn (đã có 3 định lý).
\end{itemize}

%=====================================================================
\section*{Tài liệu tham khảo}
%=====================================================================
\begin{thebibliography}{99}
\bibitem{chi2023diffusion} C. Chi et al., ``Diffusion Policy: Visuomotor Policy Learning via Action Diffusion,'' RSS 2023.
\bibitem{janner2022planning} M. Janner, Y. Du, J. B. Tenenbaum, S. Levine, ``Planning with Diffusion for Flexible Behavior Synthesis,'' ICML 2022.
\bibitem{ajay2023conditional} A. Ajay et al., ``Is Conditional Generative Modeling all you need for Decision-Making?'' ICLR 2023.
\bibitem{lipman2023flow} Y. Lipman et al., ``Flow Matching for Generative Modeling,'' ICLR 2023.
\bibitem{zhao2023aloha} T. Z. Zhao, V. Kumar, S. Levine, C. Finn, ``Learning Fine-Grained Bimanual Manipulation with Low-Cost Hardware,'' RSS 2023.
\bibitem{ivanov2019verisig} R. Ivanov et al., ``Verisig: Verifying Safety Properties of Hybrid Systems with Neural Network Controllers,'' HSCC 2019.
\bibitem{tran2020nnv} H.-D. Tran et al., ``NNV: The Neural Network Verification Tool,'' CAV 2020.
\bibitem{huang2019reachnn} C. Huang et al., ``ReachNN: Reachability Analysis of Neural-Network Controlled Systems,'' EMSOFT 2019.
\bibitem{everett2021reachability} M. Everett, G. Habibi, C. Sun, J. P. How, ``Reachability Analysis of Neural Feedback Loops,'' IEEE Access, 2021.
\bibitem{althoff2015cora} M. Althoff, ``An Introduction to CORA 2015,'' ARCH Workshop 2015.
\bibitem{chen2013flowstar} X. Chen, E. Ábrahám, S. Sankaranarayanan, ``Flow*: An Analyzer for Non-Linear Hybrid Systems,'' CAV 2013.
\bibitem{maler2004monitoring} O. Maler, D. Nickovic, ``Monitoring Temporal Properties of Continuous Signals,'' FORMATS/FTRTFT 2004.
\bibitem{donze2010robust} A. Donzé, O. Maler, ``Robust Satisfaction of Temporal Logic over Real-Valued Signals,'' FORMATS 2010.
\bibitem{fainekos2009robustness} G. E. Fainekos, G. J. Pappas, ``Robustness of Temporal Logic Specifications for Continuous-Time Signals,'' Theoretical Computer Science, 2009.
\bibitem{deshmukh2017robust} J. V. Deshmukh et al., ``Robust Online Monitoring of Signal Temporal Logic,'' Formal Methods in System Design, 2017.
\bibitem{solar2006combinatorial} A. Solar-Lezama et al., ``Combinatorial Sketching for Finite Programs,'' ASPLOS 2006.
\bibitem{clarke2000counterexample} E. Clarke, O. Grumberg, S. Jha, Y. Lu, H. Veith, ``Counterexample-Guided Abstraction Refinement,'' CAV 2000.
\bibitem{annpureddy2011staliro} Y. Annpureddy et al., ``S-TaLiRo: A Tool for Temporal Logic Falsification for Hybrid Systems,'' TACAS 2011.
\bibitem{donze2010breach} A. Donzé, ``Breach, a Toolbox for Verification and Parameter Synthesis of Hybrid Systems,'' CAV 2010.
\bibitem{dreossi2019verifai} T. Dreossi et al., ``VerifAI: A Toolkit for the Formal Design and Analysis of AI-based Systems,'' CAV 2019.
\bibitem{wabersich2021predictive} K. P. Wabersich, M. N. Zeilinger, ``A Predictive Safety Filter for Learning-Based Control of Constrained Nonlinear Dynamical Systems,'' Automatica, vol.~129, 109597, 2021.
\bibitem{bastani2021safe} O. Bastani, ``Safe Reinforcement Learning with Nonlinear Dynamics via Model Predictive Shielding,'' ACC 2021.
\bibitem{ames2017cbf} A. D. Ames, X. Xu, J. W. Grizzle, P. Tabuada, ``Control Barrier Function Based Quadratic Programs for Safety Critical Systems,'' IEEE TAC, vol.~62, no.~8, 2017.
\bibitem{ames2019cbf-ecc} A. D. Ames et al., ``Control Barrier Functions: Theory and Applications,'' ECC 2019.
\bibitem{cheng2019end} R. Cheng, G. Orosz, R. M. Murray, J. W. Burdick, ``End-to-End Safe Reinforcement Learning through Barrier Functions,'' AAAI 2019.
\bibitem{xiao2023safediffuser} W. Xiao, T.-H. Wang, C. Gan, D. Rus, ``SafeDiffuser: Safe Planning with Diffusion Probabilistic Models,'' arXiv:2306.00148, 2023; ICLR 2025.
\bibitem{power2023sampling} T. Power, R. Soltani-Zarrin, S. Iba, D. Berenson, ``Sampling Constrained Trajectories Using Composable Diffusion Models,'' IROS Workshop 2023.
\bibitem{mayne2000constrained} D. Q. Mayne, J. B. Rawlings, C. V. Rao, P. O. M. Scokaert, ``Constrained Model Predictive Control: Stability and Optimality,'' Automatica, vol.~36, no.~6, 2000.
\bibitem{yu2024modelpredictive} X. Yu, W. Dong, X. Yin, S. Li, ``Model Predictive Monitoring of Dynamical Systems for Signal Temporal Logic Specifications,'' arXiv:2209.12493; Automatica, 2024.
\bibitem{pathconsistent2025} Anonymous/authors, ``From Demonstrations to Safe Deployment: Path-Consistent Safety Filtering for Diffusion Policies,'' arXiv:2511.06385, 2025.
\end{thebibliography}

\end{document}
